% Introduction
\part{Introduction}

\chapter{Introduction}
\section{General Introduction}

Cyber threats continue to evolve in scale and sophistication, impacting private organizations as well as public institutions.
Large-scale incident analyses show that modern attacks remain effective because adversaries repeatedly exploit common
initial access vectors and operational weaknesses, leading to significant business and mission impact \cite{verizonDBIR2025}.

In this context, penetration testing is widely used to evaluate security by simulating real attack behavior under controlled
conditions. According to NIST, penetration testing verifies the extent to which a system, device, or process can resist active
attempts to compromise its security, and often involves combining multiple weaknesses to gain deeper access than any single
vulnerability would allow \cite{nistPenTest,nistSmallBizGlossary}. Methodological guidance such as NIST SP 800-115 further
emphasizes structured planning, execution, analysis of findings, and mitigation-focused reporting \cite{nistSP800115}.

At the same time, recent advances in large language models (LLMs) and agentic AI have enabled systems capable of planning and
executing multi-step technical tasks. Research has begun exploring how multi-agent LLM frameworks can emulate the collaborative
workflow of penetration testing teams to improve efficiency and reduce manual effort \cite{vulnbot}.

Therefore, this project aims to leverage a hybrid multi-agent approach to support and improve the penetration testing process,
while enforcing strict control over scope, safety, and traceability of actions and results. The goal is to increase the
repeatability and scalability of testing activities without compromising governance and operational constraints.
\lipsum[1] \\
